Large language model agents can acquire complex capabilities through multi-step interaction and tool use, but their trajectories can also be illegally collected to distill student agents. However, existing watermarking methods either do not fit the structured and interactive nature of agent environments or lack reliable effectiveness across tasks and model architectures. We introduce \name{}, a behavioral watermarking framework for tracing unauthorized agent distillation. \name{} dynamically inserts safe, non-essential auxiliary action into teacher trajectories, and stores the associated contexts as private evidence cards. To audit a suspicious student model, \name{} constructs paired real and fake probes from these cards and applies a card-level sign test. This black-box protocol supports both model-level detection and trace-level attribution. Across three agent benchmarks, two teacher agents, and four student architectures, \name{} detects all 24 distilled models with zero false positives on 48 clean models. It also preserves task utility and remains effective against data flooding, paraphrasing, truncation, and adaptive cleaning attacks. Our code will be released at \href{https://github.com/qx041609/Auxmark}{this URL}.